%% =====================================================================
%%  SmArchitect — Q1 paper (single-source, multi-format).
%%  Compile on Overleaf as-is (IEEEtran + BibTeX), or convert to any
%%  format with `python build.py <pdf|html|docx|md|all>`.
%%
%%  No external figure assets are required: figures are self-contained
%%  framed placeholders (\figplaceholder) until the experiment scripts
%%  emit real artifacts. Experimental numbers are \result{...} markers,
%%  never fabricated.
%% =====================================================================
\documentclass[journal]{IEEEtran}

\usepackage{amsmath,amssymb,amsfonts}
\usepackage{amsthm}
\usepackage{booktabs}
\usepackage{array}
\usepackage{multirow}
\usepackage{enumitem}
\usepackage{algorithm}
\usepackage{algpseudocode}
\usepackage{graphicx}
\usepackage{xcolor}
\usepackage{url}
\usepackage[hidelinks]{hyperref}

% ----- draft toggles -------------------------------------------------
\newif\ifdraft\drafttrue   % \draftfalse for camera-ready (hides markers)
\definecolor{resultblue}{RGB}{20,90,170}
\definecolor{todored}{RGB}{170,30,30}
\newcommand{\result}[1]{\ifdraft\textcolor{resultblue}{\textbf{[R:\,#1]}}\else#1\fi}
\newcommand{\todo}[1]{\ifdraft\textcolor{todored}{\textbf{[TODO:\,#1]}}\else\fi}

% ----- self-contained figure placeholder (no image files needed) -----
% \figplaceholder{height}{what it shows}{how it is generated}
\newcommand{\figplaceholder}[3]{%
  \setlength{\fboxsep}{8pt}%
  \fbox{\begin{minipage}[c][#1][c]{0.92\linewidth}\centering
    \footnotesize\itshape #2\\[4pt]
    \textcolor{resultblue}{\scriptsize generated by: #3}
  \end{minipage}}%
}

% ----- theorem environments -----------------------------------------
\theoremstyle{plain}
\newtheorem{theorem}{Theorem}
\newtheorem{proposition}{Proposition}
\theoremstyle{definition}
\newtheorem{definition}{Definition}

% ----- handy macros for the formal model ----------------------------
\newcommand{\G}{G}
\newcommand{\Gset}{\mathcal{G}}
\newcommand{\Szt}{S_{\mathrm{ZT}}}
\newcommand{\App}{\mathrm{App}}
\newcommand{\Subj}{\mathrm{Subj}}
\newcommand{\pub}{\mathrm{pub}}
\newcommand{\reach}{\rightsquigarrow}
\newcommand{\code}[1]{\texttt{\small #1}}

\begin{document}

\title{Zero Trust as a Satisfaction Function: Formal Multi-Objective
Synthesis of Auditable Cloud Architectures with Human-in-the-Loop LLM Agents}

\author{First~Author,~Second~Author,~and~Third~Author%
\thanks{Manuscript prepared for submission to IEEE Transactions on Dependable
and Secure Computing. Author affiliations and contact details omitted for
review. Artifact (the SmArchitect platform) available to reviewers.}}

\markboth{Preprint --- under review}{Author \MakeLowercase{\textit{et al.}}: Zero Trust as a Satisfaction Function}

\maketitle

\begin{abstract}
Cloud architectures are predominantly designed by intuition and audited against
\emph{qualitative} Zero-Trust checklists, leaving a gap between what is
\emph{optimizable} and what is \emph{verifiable}. We recast architecture design
as a constrained multi-objective problem over a provider-agnostic \emph{Cloud
Resource Graph} (CRG). Zero Trust becomes a formally defined \emph{satisfaction
function} $\Szt$ mapped to the controls of NIST SP~800-207/207A, and we prove
that it generalizes the scoring engine deployed in our platform. Cost,
availability, and a graph-theoretic attack surface form an explicit objective
vector whose \emph{non-dominated} (Pareto-optimal) architectures we compute
exactly via a mixed-integer linear program (MILP) swept with the
$\varepsilon$-constraint method under HIPAA/PCI hard constraints. A
human-in-the-loop pipeline of large language model (LLM) agents acts as a set of
stochastic \emph{search operators} whose proposals are projected onto the
feasible manifold by a deterministic planner---a neuro-symbolic loop that we
prove is monotone non-decreasing in $\Szt$ under remediation. We further show
that the platform's one-click remediation is \emph{minimum-cost} for per-node
controls. The framework is implemented on a working platform (SmArchitect) and
evaluated on \result{$5$} architecture families across \result{$3$} compliance
overlays, measuring Pareto hypervolume, remediation optimality gap, IaC
deployability, and stage-wise agent ablation. Results are reproducible from the
released artifact.
\end{abstract}

\begin{IEEEkeywords}
Zero Trust, cloud architecture, multi-objective optimization, mixed-integer
linear programming, neuro-symbolic AI, large language model agents,
Infrastructure-as-Code, NIST SP 800-207.
\end{IEEEkeywords}

% =====================================================================
\section{Introduction}
\IEEEPARstart{C}{loud} architects routinely make decisions---multi-AZ or
single-AZ? public load balancer or private ingress? where to terminate TLS?---
that simultaneously determine \emph{cost}, \emph{availability}, and
\emph{security posture}. Yet the artifacts that govern these decisions are
mismatched. Cost is estimated in spreadsheets; availability is reasoned about
informally; and Zero-Trust ``compliance'' is assessed against maturity-model
checklists that are qualitative and not mechanically tied to the architecture
itself. There is no single formal object that is at once \emph{optimizable} (we
can search for better designs) and \emph{verifiable} (we can prove a design
satisfies a named control).

This paper closes that gap. We treat an architecture as a \emph{Cloud Resource
Graph} (CRG): a provider-agnostic, attribute-bearing directed graph in which
nodes are cloud resources carrying security, cost, high-availability (HA), and
scaling attributes, and edges are network/data/control flows with sensitivity
flags. Over this single object we define (i) a Zero-Trust \emph{satisfaction
function} $\Szt$ grounded in NIST SP~800-207 and 800-207A; (ii) an explicit
objective vector adding cost, availability, and a graph-theoretic attack
surface; and (iii) a constrained multi-objective program whose Pareto-optimal
solutions are the ``optimal secure cloud systems.'' A pipeline of LLM agents
proposes and edits graphs; a deterministic planner projects every proposal back
onto the feasible manifold, giving a neuro-symbolic synthesis loop with formal
guarantees.

Our contributions are:
\begin{enumerate}[leftmargin=1.4em]
\item A formal CRG model with a Zero-Trust satisfaction function $\Szt$ that we
prove \emph{generalizes} the scoring engine deployed in our platform
(Theorem~\ref{thm:equiv}), plus a computable attack-surface metric $A(\G)$.
\item An exact \emph{MILP + $\varepsilon$-constraint} synthesis that returns the
Pareto frontier over cost, $\Szt$, availability, and attack surface subject to
HIPAA/PCI hard constraints.
\item A proof that one-click remediation is \emph{minimum-cost} for per-node
controls (Theorem~\ref{thm:remediation}), and that the agentic remediate-all
policy is monotone non-decreasing in $\Szt$ and convergent
(Theorem~\ref{thm:monotone}).
\item A neuro-symbolic agent pipeline---LLM proposers plus a deterministic
feasibility projector---evaluated for deployability, optimality gap, and ablated
stage-wise contribution on real templates.
\end{enumerate}

The framework is realized on \emph{SmArchitect}, a working five-plane platform;
every equation in this paper maps to a concrete module, and every reported
number is regenerated from the released artifact.

% =====================================================================
\section{Related Work}\label{sec:related}

\subsection{Cloud architecture recommendation and model-driven optimization}
Model-driven approaches such as CloudMIG~\cite{frey2011cloudmig} and topology
standards such as TOSCA~\cite{tosca2013} capture cloud deployments as models,
and a long line of work optimizes \emph{placement} (VM/container scheduling).
These operate at a different abstraction level from architecture-time
\emph{design} and do not treat Zero-Trust compliance as a first-class,
graph-verifiable objective.

\subsection{Zero-Trust maturity and policy-as-code}
NIST SP~800-207~\cite{nist800207} and its cloud-native companion
800-207A~\cite{nist800207a} define Zero-Trust tenets and enforcement components
(policy enforcement points, workload identity). In practice these are operating
as qualitative maturity checklists or as runtime policy-as-code (OPA/Rego
\cite{opa}) gating, with machine-readable catalogs in OSCAL~\cite{oscal}. None
provides a continuous, optimizable satisfaction function over an architecture
graph, nor a minimum-cost remediation operator.

\subsection{LLM-for-IaC and agentic synthesis}
Reasoning-and-acting agents~\cite{yao2023react} and emerging LLM-for-IaC tools
generate Terraform/Pulumi from natural language~\cite{llmiac2024}, but lack a
formal compliance/cost objective or optimality guarantees and do not project
proposals onto a verified feasible set. We frame the LLM as a stochastic
proposal operator within a neuro-symbolic loop~\cite{garcez2023neurosymbolic},
which yields the guarantees these tools lack.

Table~\ref{tab:related} positions this work: the combination of a
\emph{formal} Zero-Trust objective, \emph{exact} multi-objective optimization,
\emph{deployable} IaC, and a \emph{human-in-the-loop} agentic loop is, to our
knowledge, unique.

% =====================================================================
\section{Background: The SmArchitect Platform}\label{sec:background}
SmArchitect organizes the system into five planes over one CRG: a \emph{Canvas}
plane (interactive editor), an \emph{Agent} plane (a staged LLM pipeline:
Requirements $\to$ Research $\to$ Architect $\to$ Cost $\to$ Security $\to$
Compliance $\to$ IaC with human checkpoints), a \emph{Graph} plane (the CRG as
single source of truth, versioned), a \emph{Policy} plane (Zero-Trust controls
ZT-1\,$\ldots$\,ZT-8 plus HIPAA/SOC2/PCI overlays), and an \emph{IaC} plane
(export to Pulumi/Terraform/Helm/CloudFormation and import with drift
detection). Figure~\ref{fig:planes} shows the planes and the modules realizing
them. The Policy plane's eight controls are aligned to the seven tenets of NIST
800-207 and the cloud-native enforcement model of 800-207A; their precise
predicates are given in Section~\ref{sec:formal} and Table~\ref{tab:controls}.

\begin{figure}[t]\centering
\figplaceholder{3.2cm}{Five-plane SmArchitect architecture: Canvas, Agent,
Graph, Policy, IaC, annotated with the modules that realize each plane and the
prompt$\to$spec$\to$CRG$\to$optimizer$\to$remediated-CRG$\to$IaC data flow.}%
{redraw from SmArchitect\_Platform\_Plan.md \S5}
\caption{System architecture (five planes over one CRG).}
\label{fig:planes}
\end{figure}

% =====================================================================
\section{Formal Model}\label{sec:formal}
Table~\ref{tab:notation} summarizes notation.

\subsection{Cloud Resource Graph}
\begin{definition}[CRG]\label{def:crg}
A Cloud Resource Graph is a tuple $\G=(V,E,\Phi_V,\Phi_E)$ where $V$ is a set of
resource nodes, $E\subseteq V\times V$ is a set of directed edges, and
\[
\Phi_V(v)=\big(\mathrm{prov},\mathrm{svc},\mathrm{cat},\mathrm{tier},\mathrm{reg},h(v),\sigma(v),sc(v),\kappa(v)\big).
\]
The security vector $\sigma(v)\in\{0,1\}^7$ encodes
$(\text{auth},\text{encT},\text{encR},\text{pub},\text{mtls},\text{waf},\text{priv})$,
mapping one-to-one to the platform's \code{SecuritySettings}. The HA mode
$h(v)\in\{\textsf{single},\textsf{az},\textsf{region},\textsf{global}\}$, and
$sc(v)$, $\kappa(v)$ are scaling and cost attributes. Each edge carries
$\Phi_E(e)=(\text{kind}\in\{\text{net},\text{data},\text{ctrl}\},\text{proto},\text{latCrit},s(e))$
with sensitivity $s(e)\in\{0,1\}$.
\end{definition}

We add a virtual internet source $v_\infty$ with $v_\infty\to v$ iff
$\pub(v)=1$, and a trust boundary
\begin{equation}
\tau(v)=\begin{cases}
\textsf{trusted} & \neg\pub(v)\wedge\text{priv}(v)\\
\textsf{dmz} & \mathrm{tier}(v)\in\{\text{edge},\text{ingress}\}\wedge\pub(v)\\
\textsf{untrusted} & \text{otherwise.}
\end{cases}
\label{eq:tau}
\end{equation}

\begin{definition}[Attack surface]\label{def:attack}
Let $D=\{v:\mathrm{cat}(v)\in\{\text{data},\text{storage}\}\}$ and let
$v_\infty\reach d$ denote directed reachability. Then
\begin{equation}
A(\G)=\sum_{v\in V}\pub(v)\,\omega_{\mathrm{tier}(v)}
      +\sum_{d\in D}\mathbb{1}[v_\infty\reach d]\big(1+\!\!\max_{e\in\delta(d)}\! s(e)\big),
\label{eq:attack}
\end{equation}
with tier exposure weights $\omega$. $A(\G)$ is computed by breadth-first
reachability from $v_\infty$ (new module \code{security/attack-surface.ts} /
\code{server/app/security.py}).
\end{definition}

\subsection{Zero-Trust Satisfaction Function}
Each control $c_i$ ($i=1,\ldots,8$) has an applicable subject set
$\Subj_i(\G)\subseteq V$ and a per-subject predicate $p_i$:
\begin{equation}
c_i(\G)=\begin{cases}
\textsf{n/a} & \Subj_i=\varnothing\\
1 & \forall v\in\Subj_i:\,p_i(v)\\
0 & \forall v\in\Subj_i:\,\neg p_i(v)\\
0.5 & \text{otherwise.}
\end{cases}
\label{eq:ci}
\end{equation}
Controls $c_1,c_2,c_4,c_5,c_8$ are \emph{per-node} (predicate over $\sigma(v)$);
$c_3$ (policy enforcement point in front of public compute), $c_6$
(observability present), and $c_7$ (central secret store present) are
\emph{structural}. Table~\ref{tab:controls} gives each $(\Subj_i,p_i)$ and its
NIST reference. Let $\App(\G)=\{i:c_i\neq\textsf{n/a}\}$. The deployed engine
computes the uniform score
\begin{equation}
S_{\text{uni}}(\G)=\frac{100}{|\App|}\sum_{i\in\App}c_i(\G),
\label{eq:suni}
\end{equation}
which we generalize with NIST-tenet criticality weights $w_i$:
\begin{equation}
\Szt(\G)=100\cdot\frac{\sum_{i\in\App}w_i\,c_i(\G)}{\sum_{i\in\App}w_i}.
\label{eq:szt}
\end{equation}

\begin{theorem}[Equivalence to the deployed engine]\label{thm:equiv}
The platform routine \code{evaluateFramework()} computes
$\big\lfloor S_{\text{uni}}(\G)\big\rceil$. Hence $\Szt$ with $w_i\equiv 1$
equals the deployed Zero-Trust score (up to rounding).
\end{theorem}
\begin{proof}[Proof sketch]
The implementation accumulates
$\texttt{earned}=\sum_{i\in\App}c_i(\G)$ with the same pass/partial/fail values
$\{1,0.5,0\}$ as \eqref{eq:ci}, sets $|\App|=\texttt{applicable.length}$, and
returns $\mathrm{round}(\texttt{earned}/|\App|\cdot100)$, which is exactly
$\lfloor S_{\text{uni}}\rceil$. Setting $w_i\equiv1$ in \eqref{eq:szt} reduces it
to \eqref{eq:suni}.
\end{proof}

\subsection{Remediation as Minimum-Cost Graph Edit}
Let $d(\G,\G')=\sum_{\text{flips}}\theta_{\text{attr}}+\sum_{\text{adds}}\theta_{\text{node}}$
be a weighted graph-edit distance. A remediating transform for control $c_i$ is
\begin{equation}
R_{c_i}(\G)=\arg\min_{\G':\,c_i(\G')=1} d(\G,\G').
\label{eq:remediate}
\end{equation}
\begin{theorem}[Minimal remediation]\label{thm:remediation}
For the per-node controls $c_1,c_2,c_4,c_5,c_8$, the platform's \code{remediate()}
flips exactly the violating security bits and is the unique $d$-minimizer. For
the structural controls $c_3,c_6,c_7$ it activates exactly one catalog node and
is minimal under the augmentation set.
\end{theorem}
\begin{proof}[Proof sketch]
A per-node predicate $p_i$ depends only on $\sigma(v)$; achieving $c_i(\G')=1$
requires $p_i(v)$ for every violating $v\in\Subj_i$, and each such $v$ needs at
least one bit flip. The implementation flips precisely those bits and no others,
so its edit set is contained in every feasible edit set and has minimal
$\theta_{\text{attr}}$ weight. Structural controls are satisfied by the presence
of one enforcement/observability/secret node; adding a single least-cost catalog
node is therefore minimal, since the empty edit leaves $c_i=0$.
\end{proof}

\subsection{Multi-Objective Synthesis}
We optimize over decision variables on a fixed (optionally augmentable)
topology: per-node security bits $\sigma_v\in\{0,1\}^7$, HA one-hot $y_{v,m}$
($\sum_m y_{v,m}=1$), provider one-hot $\pi_{v,p}$ for provider-flexible nodes,
and augmentation-node activations $a_u\in\{0,1\}$. The objective vector is
\begin{equation}
F(\G)=\big(C(\G),\,-\Szt(\G),\,-R(\G),\,A(\G)\big),
\label{eq:F}
\end{equation}
with
\begin{align}
C(\G)&=\sum_{v}\mathrm{base}(v,\pi)\sum_{m} f_m\,y_{v,m}\,\mathrm{units}(v),\quad f\in\{1,1.6,2.4,3\},\label{eq:cost}\\
R(\G)&=\sum_{v}\rho(h_v)\,\mathrm{crit}(v),\quad \rho=(0.5,0.8,0.95,1.0).\label{eq:resil}
\end{align}
Equation~\eqref{eq:cost} reproduces the deployed cost model
($\mathrm{base}\times\text{HA-factor}\times\text{units}$). The hard constraints
are budget $C(\G)\le B$; HIPAA $\Rightarrow \text{encR}(d)=\text{encT}(d)=1$ for
data nodes; PCI $\Rightarrow \pub(d)=0$ and $\text{waf}(\text{edge})=1$;
structural validity (an ingress$\to$compute$\to$data path exists, encoded as a
flow constraint); and deployability $\mathrm{svc}(v)\in\mathrm{TF\_TYPE}$.

\paragraph{MILP.}
Per-node controls are linear in $\sigma$. Structural controls are encoded with
big-$M$/indicator constraints over augmentation activations $a_u$ and adjacency
(e.g.\ $c_3$: a public compute node requires an active authenticating ingress on
its in-path). Bilinear cost terms are linearized by treating per-node units as
constants and selecting provider via $\pi$. The result is a MILP solved with
PuLP/CBC.

\paragraph{$\varepsilon$-constraint Pareto.}
We obtain the exact frontier by scalarizing to a single objective and bounding
the rest:
\begin{equation}
\min_{\G}\ C(\G)\quad\text{s.t.}\quad \Szt(\G)\ge\varepsilon_S,\ R(\G)\ge\varepsilon_R,\ A(\G)\le\varepsilon_A,
\label{eq:eps}
\end{equation}
sweeping the $(\varepsilon_S,\varepsilon_R,\varepsilon_A)$ grid. Every optimum of
\eqref{eq:eps} is Pareto-optimal, and the sweep recovers the non-dominated set
\begin{equation}
\mathcal{P}^*=\{\G:\nexists\,\G'\ \text{feasible},\ F(\G')\preceq F(\G),\ F(\G')\neq F(\G)\}.
\label{eq:pareto}
\end{equation}
An ``optimal secure cloud system'' is any member of $\mathcal{P}^*$; we report
frontier quality with the hypervolume indicator $HV(\mathcal{P}^*;r)$.

\subsection{Agentic Pipeline as Formal Search}
A search state is $s=(\G,\text{spec},\text{stage})$ with actions
$\mathcal{A}=\{\textsf{gen\_template},\textsf{apply\_patch},\textsf{remediate\_control},\textsf{recompute\_cost},\textsf{switch\_provider}\}$.
An LLM agent is a stochastic proposal operator $\omega:s\mapsto\Delta(\mathcal{A})$;
the deterministic planner is a feasibility projector
$\Pi=\textsf{validateGraph}\circ\textsf{recomputeCosts}\circ\textsf{autoLayout}$,
which is idempotent ($\Pi\circ\Pi=\Pi$) and repairs invalid proposals (dangling
edges, hallucinated resources). A human checkpoint is an approval filter on state
transitions.

\begin{theorem}[Monotone convergence]\label{thm:monotone}
Under the remediate-all policy $\rho_{\mathrm{RA}}=\bigcirc_i R_{c_i}$,
$\Szt(\G_{t+1})\ge\Szt(\G_t)$ for all $t$, and the sequence converges to a
feasible $\G^\star$ with $\Szt(\G^\star)=100$ whenever a remediation exists for
every applicable control.
\end{theorem}
\begin{proof}[Proof sketch]
Each $R_{c_i}$ either flips security bits toward satisfaction or adds a passing
node; both operations are monotone in every $c_j$ (added security or added
passing subjects never make a satisfied predicate fail). Hence $\Szt$ is
non-decreasing. As $\Szt$ is bounded above by $100$ and $\App$ is finite, the
process terminates; if every applicable control admits a remediation (true for
$c_1,\ldots,c_8$ by construction), all reach $c_i=1$ and $\Szt=100$. The
projector $\Pi$ preserves structural validity throughout.
\end{proof}

% =====================================================================
\section{Algorithms}\label{sec:algorithms}

\begin{algorithm}[t]
\caption{$\textsc{EvaluateZT}(\G)$ --- Zero-Trust satisfaction}
\label{alg:eval}
\begin{algorithmic}[1]
\State $\App\gets\varnothing;\ \texttt{earned}\gets0$
\For{each control $c_i$, $i=1\ldots8$}
  \State $S\gets\Subj_i(\G)$
  \If{$S=\varnothing$} \State \textbf{continue} \Comment{n/a} \EndIf
  \State $\App\gets\App\cup\{i\}$
  \State $P\gets\{v\in S:p_i(v)\}$
  \State $c_i\gets \mathbf{1}[P{=}S]\cdot1 + \mathbf{1}[P{=}\varnothing]\cdot0 + \mathbf{1}[\varnothing\subsetneq P\subsetneq S]\cdot0.5$
  \State $\texttt{earned}\gets\texttt{earned}+w_i\,c_i$
\EndFor
\State \Return $100\cdot\texttt{earned}/\textstyle\sum_{i\in\App}w_i$
\end{algorithmic}
\end{algorithm}

\begin{algorithm}[t]
\caption{$\textsc{ParetoSweep}(\G,B,\text{tags})$ --- exact frontier}
\label{alg:pareto}
\begin{algorithmic}[1]
\State $M\gets\textsc{BuildMILP}(\G,B,\text{tags})$ \Comment{vars $\sigma,y,\pi,a$; Eqs.\,\eqref{eq:cost}--\eqref{eq:eps}}
\State $\mathcal{P}\gets\varnothing$
\For{each grid point $(\varepsilon_S,\varepsilon_R,\varepsilon_A)$}
  \State add constraints $\Szt\ge\varepsilon_S,\ R\ge\varepsilon_R,\ A\le\varepsilon_A$ to $M$
  \State $\G^\*\gets\arg\min C(\G)$ via CBC; \textbf{if} feasible \textbf{then} $\mathcal{P}\gets\mathcal{P}\cup\{\G^\*\}$
  \State remove the $\varepsilon$ constraints
\EndFor
\State \Return non-dominated filter of $\mathcal{P}$ \Comment{Eq.\,\eqref{eq:pareto}}
\end{algorithmic}
\end{algorithm}

\begin{algorithm}[t]
\caption{$\textsc{AgentSynthesize}(\text{spec})$ --- neuro-symbolic loop}
\label{alg:agent}
\begin{algorithmic}[1]
\State $\G\gets\Pi(\textsc{GenTemplate}(\text{spec}))$
\For{stage in $[\,$Research, Architect, Cost, Security, Compliance, IaC$\,]$}
  \State $a\sim\omega_{\text{stage}}(\G,\text{spec})$ \Comment{LLM proposal}
  \State $\G'\gets\Pi(\text{apply}(a,\G))$ \Comment{project to feasible}
  \If{$\textsc{HumanApproves}(\G,\G')$} $\G\gets\G'$ \EndIf
\EndFor
\State \Return $\rho_{\mathrm{RA}}(\G)$ \Comment{Theorem~\ref{thm:monotone}}
\end{algorithmic}
\end{algorithm}

% =====================================================================
\section{Implementation}\label{sec:impl}
The framework is implemented on SmArchitect (TypeScript canvas/engines;
Python/FastAPI agent, policy, and IaC services; LiteLLM gateway). The CRG type
system, the Zero-Trust evaluation engine \eqref{eq:ci}--\eqref{eq:suni}, the cost
model \eqref{eq:cost}, the deterministic planner ($\Pi$), the agent pipeline, and
the IaC exporters already exist. The formal contributions add: the attack-surface
module \eqref{eq:attack}; the weighted score \eqref{eq:szt}; unification of the
backend control set to all eight controls (it currently implements five); and the
MILP/$\varepsilon$-constraint optimizer and experiment harness. Table~\ref{tab:roadmap}
maps each component to its module.

% =====================================================================
\section{Evaluation}\label{sec:eval}
\emph{All numbers below are produced by the released harness; placeholders mark
quantities to be filled from \code{results/}.}

\paragraph{Datasets} Five architecture families generated from prompts
(three-tier, serverless, microservices, data pipeline, multi-region DR), three
compliance overlays (Zero Trust, HIPAA, PCI), and one brownfield graph imported
from real Terraform state. Table~\ref{tab:exp} lists configurations.

\paragraph{Metrics} (a)~$\Szt$ before/after remediation; (b)~cost delta
$\Delta C$ of remediation; (c)~Pareto hypervolume; (d)~optimality gap of
heuristic vs.\ MILP-optimal remediation; (e)~agent time-to-feasible and number of
human checkpoints; (f)~IaC deployability (\% generated Terraform passing
\code{terraform validate}); (g)~talk-to-edit success on 20 canonical commands;
(h)~Zero-Trust precision/recall vs.\ a manual control audit.

\paragraph{Baselines} Manual-architect (expert), template-only (no
optimization), LLM-only (no planner, no formal Zero Trust), and heuristic vs.\
MILP-optimal remediation.

\paragraph{Statistical protocol} Stochastic (LLM) configurations are run for
$\ge 30$ trials reporting mean $\pm$ 95\% CI; methods are compared with paired
Wilcoxon signed-rank tests under Holm correction; model and version are logged
per run via LiteLLM, sweeping $\ge 2$ models (frontier and cheap) for the
cost/latency trade-off.

\paragraph{Results (to populate)}
The MILP recovers a frontier of \result{$n$} non-dominated architectures with
hypervolume \result{$HV$}; automated remediation raises mean $\Szt$ from
\result{$x$} to \result{$100$} at cost delta \result{$\Delta C$}; heuristic
remediation matches the MILP optimum on per-node controls with an optimality gap
of \result{$g\%$} on structural controls; \result{$p\%$} of generated Terraform
passes \code{terraform validate}; ablating the Security Critic and Compliance
stages reduces $\Szt$ by \result{$\Delta\Szt$}.
Figures~\ref{fig:zt}--\ref{fig:provider} and Table~\ref{tab:exp} report the full
results.

\begin{figure}[t]\centering
\figplaceholder{2.7cm}{ZT-1\,$\ldots$\,ZT-8 evaluation on a sample three-tier
graph: each node colored pass/partial/fail.}{compliance engine + studio lens}
\caption{Zero-Trust control evaluation flow.}
\label{fig:zt}
\end{figure}

\begin{figure}[t]\centering
\figplaceholder{3.0cm}{Pareto frontier (cost vs.\ $\Szt$) over $\ge20$
configurations of one template; attack surface as color, availability as marker
size; non-dominated front highlighted.}{server/app/optimizer.py $\to$ CSV $\to$ matplotlib}
\caption{Exact multi-objective Pareto frontier.}
\label{fig:pareto}
\end{figure}

\begin{figure}[t]\centering
\figplaceholder{2.4cm}{Before/after remediation diff: node statuses flip to pass
and the cost delta is annotated.}{compliance engine + diff.ts}
\caption{Remediation graph diff.}
\label{fig:remediate}
\end{figure}

\begin{figure}[t]\centering
\figplaceholder{2.4cm}{Provider comparison (AWS/Azure/GCP/private) per template
from \code{compareProviders()}.}{cost.ts}
\caption{Multi-cloud cost comparison.}
\label{fig:provider}
\end{figure}

% =====================================================================
\section{Discussion: Limitations and Threats to Validity}\label{sec:discussion}
The cost model is catalog-based (relative, architecture-time) rather than a live
pricing API; we report it as a comparative signal and leave live pricing to
future work. The deterministic planner's template scope bounds the generative
diversity of the agent pipeline. LLM non-determinism is mitigated by treating the
LLM as an ablated proposal operator over a deterministic feasible set, with
$\ge30$-run confidence intervals. The attack-surface metric is structural and
architecture-time; runtime exploitability and authentication of human principals
are out of scope by design. Finally, the MILP encodes structural controls via
big-$M$/indicator constraints; we validate exactness against brute force on small
templates and report any gap.

% =====================================================================
\section{Conclusion}\label{sec:conclusion}
We recast cloud-architecture design as a constrained multi-objective problem over
a unified Cloud Resource Graph, with Zero Trust as a formal satisfaction function
that provably generalizes a deployed engine, an exact MILP/$\varepsilon$-constraint
Pareto synthesis, and a neuro-symbolic agent loop with monotone remediation
guarantees. The result is a single graph that can be synthesized, verified, and
deployed. Future work includes live-pricing objectives and drift-aware
re-optimization of imported estates.

% =====================================================================
% Wide tables (full column width)
% =====================================================================
\begin{table}[t]
\caption{Notation.}
\label{tab:notation}
\centering\footnotesize
\begin{tabular}{@{}ll@{}}
\toprule
Symbol & Meaning\\
\midrule
$\G=(V,E,\Phi_V,\Phi_E)$ & Cloud Resource Graph\\
$\sigma(v)\in\{0,1\}^7$ & node security vector\\
$h(v)$ & HA mode (single/az/region/global)\\
$s(e)\in\{0,1\}$ & edge data-sensitivity flag\\
$v_\infty$, $v_\infty\reach d$ & internet source; reachability\\
$\tau(v)$ & trust boundary, Eq.\,\eqref{eq:tau}\\
$A(\G)$ & attack surface, Eq.\,\eqref{eq:attack}\\
$c_i(\G)\in\{0,0.5,1\}$ & control satisfaction, Eq.\,\eqref{eq:ci}\\
$\App(\G)$ & applicable controls\\
$\Szt(\G)$ & Zero-Trust score, Eq.\,\eqref{eq:szt}\\
$w_i$ & control weight (NIST tenet criticality)\\
$R_{c_i}$ & remediating transform, Eq.\,\eqref{eq:remediate}\\
$F(\G)$ & objective vector, Eq.\,\eqref{eq:F}\\
$C,R,A$ & cost, resilience, attack surface\\
$\mathcal{P}^*$ & Pareto set, Eq.\,\eqref{eq:pareto}\\
$\Pi$ & feasibility projector (planner)\\
$\omega$ & LLM proposal operator\\
\bottomrule
\end{tabular}
\end{table}

\begin{table*}[t]
\caption{Positioning against related work. \result{verify each cell during writing.}}
\label{tab:related}
\centering\footnotesize
\begin{tabular}{@{}lccccc@{}}
\toprule
Approach & Formal ZT objective & Multi-objective & Exact Pareto & Deployable IaC & Human-in-loop\\
\midrule
TOSCA / CloudMIG~\cite{tosca2013,frey2011cloudmig} & no & partial & no & partial & no\\
VM/container placement & no & yes & sometimes & no & no\\
ZT maturity models~\cite{nist800207} & qualitative & no & no & no & no\\
OPA/Rego policy-as-code~\cite{opa} & boolean gate & no & no & runtime & no\\
LLM-for-IaC~\cite{llmiac2024} & no & no & no & yes & partial\\
\textbf{This work} & \textbf{yes ($\Szt$)} & \textbf{yes} & \textbf{yes} & \textbf{yes} & \textbf{yes}\\
\bottomrule
\end{tabular}
\end{table*}

\begin{table*}[t]
\caption{Zero-Trust control to CRG-field mapping (ZT-1\,$\ldots$\,ZT-8).}
\label{tab:controls}
\centering\footnotesize
\begin{tabular}{@{}llllc@{}}
\toprule
ID & Title & Subject set $\Subj_i$ & Predicate $p_i$ / condition & NIST ref\\
\midrule
$c_1$ & Authenticate all resources & compute $\cup$ data & $\text{auth}(v)$ & 800-207 \S2.1 T1\\
$c_2$ & Encrypt in transit & non-edge $\cup$ compute & $\text{encT}(v)$ & 800-207 \S2.1 T2\\
$c_3$ & Per-session via PEP & public compute & authenticating ingress on in-path & 800-207 \S3.1\\
$c_4$ & No implicit network trust & data $\cup$ storage & $\neg\pub(v)\wedge\text{priv}(v)$ & 800-207 \S2.1 T6\\
$c_5$ & Workload identity (mTLS) & compute w/ internal nbr & $\text{mtls}(v)$ & 800-207A \S3\\
$c_6$ & Continuous monitoring & graph-global & observability node present & 800-207 \S2.1 T7\\
$c_7$ & Central secret store & data present & secret-manager node present & 800-207 \S3.4.1\\
$c_8$ & WAF at edge & public edge/ingress & $\text{waf}(v)$ & 800-207A \S4\\
\bottomrule
\end{tabular}
\end{table*}

\begin{table*}[t]
\caption{Experiment configurations and result placeholders.}
\label{tab:exp}
\centering\footnotesize
\begin{tabular}{@{}llllcccc@{}}
\toprule
Template & Overlay & Budget $B$ & Models & $\Szt$ pre & $\Szt$ post & $\Delta C$ & TF-valid \%\\
\midrule
three-tier & ZT & \$1.5k & opus/haiku & \result{--} & \result{--} & \result{--} & \result{--}\\
serverless & HIPAA & \$1k & opus/haiku & \result{--} & \result{--} & \result{--} & \result{--}\\
microservices & ZT+PCI & \$3k & opus/haiku & \result{--} & \result{--} & \result{--} & \result{--}\\
data-pipeline & ZT & \$4k & opus/haiku & \result{--} & \result{--} & \result{--} & \result{--}\\
multi-region DR & ZT+HIPAA & \$6k & opus/haiku & \result{--} & \result{--} & \result{--} & \result{--}\\
\bottomrule
\end{tabular}
\end{table*}

\begin{table}[t]
\caption{Implementation roadmap (module mapping).}
\label{tab:roadmap}
\centering\footnotesize
\begin{tabular}{@{}lll@{}}
\toprule
Component & New/Extend & Module\\
\midrule
$\Szt$ weighted & extend & \code{compliance/shared.ts}, \code{compliance.py}\\
$A(\G)$ & new & \code{security/attack-surface.ts}, \code{security.py}\\
MILP + Pareto & new & \code{server/app/optimizer.py}\\
Remediation opt. & extend & \code{compliance/zero-trust.ts}\\
Experiment runner & new & \code{server/app/experiments/}\\
Results export & new & \code{results/*.csv,*.json}\\
\bottomrule
\end{tabular}
\end{table}

% =====================================================================
\section*{Reproducibility}
Every figure and table is regenerated by \code{experiments/run.py} followed by
\code{experiments/plots.py}; deployability is checked with \code{terraform
validate}; theorem-supporting equalities are asserted by the unit tests in
\code{server/app/tests}. No number in this paper is hand-entered.

\bibliographystyle{IEEEtran}
\bibliography{references}

\end{document}
